\section{Discussion}
\label{sec:discussion}

\paragraph{Cost is a first-class objective} \ifconf
In the Windows pilot wasted tool calls cost almost as much reward as
imperfect answers, and on the full benchmark about 0.06 of reward; accuracy-only evaluation hides this. Every
lookup has latency and, with metered APIs, a price, so budgets belong in the agent and efficiency in the
benchmark. Efficiency instructions are not free: the method prompt's rule not to look up the host named in the
alert saved calls but withheld the identifier the response had to name -- a trade-off visible only in per-call
traces, not aggregate scores.
\else
Accuracy alone is a misleading measure of an investigation agent. In the Windows pilot the reward lost to wasted
tool calls was almost as large as the reward lost to imperfect answers, and on the full benchmark wasted calls
still cost about $0.06$ of reward -- a penalty that an accuracy-only evaluation hides completely. The reason is
operational: in a real SOC every lookup carries latency, and on a metered model API every lookup also carries a
price, so the discipline of spending belongs inside the agent and the measurement of that discipline belongs
inside the benchmark. Efficiency, however, is not free to ask for. The method prompt's instruction not to re-fetch
the host already named in the alert saved calls, but by the same stroke it withheld the one identifier the final
report was required to name. Trade-offs of this kind never surface in an aggregate score; they are visible only
when one reads the per-call trace of a single task.
\fi

\paragraph{Prompted method beats mechanical guards, for a strong model} \ifconf
The cache, validation, grounding
check and budgets need no model change and apply unchanged to any model behind the driver interface. We expected
them to matter; with Claude Opus 5.5 none moved reward measurably, because the guards had little to act on -- in
63 tasks the full agent repeated no call, sent five malformed ones and invented one identifier, and never reached
the hard cap. What moved efficiency was the investigation method in the prompt (start from the trigger, keep
competing hypotheses, stop when settled), cutting calls by about 30\% at unchanged quality. The order of work for
practitioners is therefore method first, guards as insurance: the guards cost nothing measurable, and their value
lies in cases this study rarely produced -- weaker models, adversarial telemetry, and runtimes that do not
validate inputs (the one malformed call that reached the v1.0 runtime crashed it, \S\ref{sec:benchmark}) -- which
\bench does not yet provoke.
\else
The cache, schema validation, grounding check and budgets require no change to the model and apply unchanged to
anything behind the driver interface, so we expected them to carry much of the agent's quality and efficiency.
With Claude Opus 5.5 they did not: none moved reward by a measurable amount, simply because they had almost nothing
to act on. Across 63 tasks the full agent repeated no call, issued only five malformed ones, invented a single
identifier, and never approached the hard cap. What actually made the agent efficient was the investigation method
written into the prompt -- start from the triggering event, hold competing hypotheses open, and stop once the
evidence settles them -- which cut tool calls by roughly 30\% at unchanged quality. For a practitioner the order of
work therefore inverts the usual intuition: write down the analyst's method first, and add the mechanical guards as
insurance. The guards remain worth keeping, because they cost nothing measurable and because their value lies
precisely in the situations this study rarely produced -- weaker or cheaper models, adversarial text in the
telemetry, and runtimes that do not validate their own inputs (the single malformed call that reached the v1.0
runtime crashed it, \S\ref{sec:benchmark}). A benchmark that provokes those situations is future work.
\fi

\ifconf\else
\paragraph{Variance is concentrated in ambiguous tasks} Run-to-run variation hardly affected tasks with a clear
answer, but one ambiguous task flipped between verdicts and alone produced reward differences of the size one
might attribute to a component. Ablations of LLM agents therefore need repeats, paired per-task comparisons and
a check of which tasks carry a difference; mean-reward tables alone credited three safeguards with harm they did
not do.
\fi

\ifconf\else
\paragraph{Guards must check the failure that actually occurs} The most frequent identifier error was not an
invented id but a real value of the wrong type -- a hostname where the host's entity id belonged. A grounding
check that accepts any string in the evidence cannot catch it; one that accepts only values seen in identifier
fields would have caught all of them at no cost. A guard should be derived from the interface's types, not a
list of anticipated mistakes, and its effect measured on traces.
\fi


\ifconf\else
\paragraph{Benchmarks built from attack recordings need cue control} Public recordings are made to be found:
they carry dataset names, lab host names and operator-chosen file names that announce an attack, so an agent
that reads them can be right for the wrong reason -- the shortcut-learning failure known from supervised
models~\cite{geirhos2020shortcut,arp2022dosdonts}, now in an agent that reads text. Our variants separate two
cues with opposite results: the agent quoted the dataset name but did not need it, yet did need the operator's
names (on the nine attack tasks whose reports cited such names it cleared 6 of 27 runs once renamed).
Evaluations on public recordings should report a mimicry condition; ours costs a string substitution and leaves
answers unchanged.
\fi

\paragraph{Label cues cut both ways} The two models disagreed about what a recorded test means. Claude Opus
5.5 read Atomic Red Team artefacts as attack evidence, GPT-6 astra as proof of an authorised exercise, and each
was consistent across repeats. Neither reading is supported by the telemetry alone, which records no
authorisation. Benchmarks built from technique tests therefore measure, in part, how a model resolves this
ambiguity, and their cue-controlled variants are needed to separate that from investigation skill.

% Shortened in the conference build (page limit); full text in the journal build. The attacker-controlled-names
% takeaway is carried by RQ7 (\S\ref{sec:remedy}) and the conclusion in the conference build.
\ifconf\else
\paragraph{Names are attacker-controlled input} The renamings that evaded the agent are those an intruder makes
anyway, and the behaviour they hid is dual-use: creating an account, building and loading a kernel module,
running a program with a preloaded library. A defender's agent should treat every name in telemetry as
attacker-controlled, justify a benign verdict by behaviour and context (who normally does this on this host,
whether a change record exists, whether the module is signed) rather than by the absence of alarming words, and
be evaluated against renaming as a matter of course. The same channel carries a stronger threat that we did not
test: text written for the agent itself~\cite{greshake2023indirect,debenedetti2024agentdojo}.

\paragraph{Benign tasks deserve equal weight} Clearing benign activity is most of the work in real SOCs. A
third of \bench is benign, and the twin tasks make it impossible to score well by keying on alert text. An
agent that over-reports would show a high false-positive rate on exactly the tasks where analysts most need help.

\paragraph{Contracts and model policies can conflict} A runtime field named \texttt{reasoning}, described as a
trace of the model's thinking, was enough to have every request declined by the provider's policy
(Appendix~\ref{app:engineering}); agent builders should treat rationale fields as statements of action and read
model reasoning from the channels the provider offers. Throughout, the agent's calls, their stated purposes and
its submissions are logged for a human reviewer, a prerequisite for assisting rather than replacing an analyst.
\fi
